\section{Interatomic Potential}

\subsection{Physical basis of metallic interatomic potentials}

Atomistic simulations require an interatomic potential to approximate the 
potential energy of a system as a function of its atomic coordinates and 
chemical identities. Given atomic positions $\{\mathbf{r}_i\}$ and species 
$\{\alpha_i\}$, the potential defines an energy landscape

\begin{equation}
    E = E(\{\mathbf{r}_i\},\{\alpha_i\}),
\end{equation}

from which atomic forces are obtained as

\begin{equation}
    \mathbf{F}_i = -\nabla_{\mathbf{r}_i} E.
\end{equation}

Unlike density functional theory (DFT), which treats 
the electronic structure explicitly within an approximate quantum-mechanical 
framework, an empirical interatomic potential represents the effects of 
electronic bonding through parameterised mathematical functions. This 
substantially reduces computational cost and enables simulations containing 
many more atoms over longer timescales.

A simple pair potential expresses the total energy as a sum of interactions 
between pairs of atoms:

\begin{equation}
    E_{\mathrm{pair}}
    = \frac{1}{2}\sum_i\sum_{j\ne i}
    \phi_{\alpha_i\alpha_j}(r_{ij}),
    \label{eq:pair-potential}
\end{equation}

where $r_{ij}=|\mathbf{r}_j-\mathbf{r}_i|$ is the interatomic separation and 
$\phi_{\alpha_i\alpha_j}$ is the pair interaction between species $\alpha_i$ 
and $\alpha_j$. Although computationally efficient, this formulation does not explicitly 
represent the dependence of an atom's bonding energy on the collective 
environment of its neighbours.

This limitation is relevant to metallic systems. An atom in a bulk metallic crystal
consequently experiences a different bonding environment from an atom at a free 
surface, near a vacancy, or within a chemically heterogeneous local 
environment. A transferable potential for such systems must therefore 
represent environmental effects rather than relying exclusively on independent 
pair interactions.

\subsection{Embedded-atom method}

The embedded-atom method (EAM) addresses this limitation by introducing an 
environment-dependent embedding energy. In its conventional form, the total 
energy is written as

\begin{equation}
    E_{\mathrm{EAM}}
    = \sum_i \left[
        F_{\alpha_i}(\rho_i)
        + \frac{1}{2}\sum_{j\ne i}
        \phi_{\alpha_i\alpha_j}(r_{ij})
    \right],
    \label{eq:eam-energy}
\end{equation}

where $F_{\alpha_i}$ is the embedding-energy function for species $\alpha_i$ 
and $\rho_i$ is the host electron density at atom $i$ contributed by neighbouring 
atoms. A typical density construction is

\begin{equation}
    \rho_i =
    \sum_{j\ne i}
    f_{\alpha_j}(r_{ij}),
    \label{eq:eam-density}
\end{equation}

where $f_{\alpha_j}$ represents the density contribution associated with a 
neighbouring atom of species $\alpha_j$. The embedding function describes the 
energy associated with placing an atom into the effective host density generated 
by its surroundings.

However, conventional EAM represents the local environment primarily through a scalar 
density. Two environments with similar effective densities but different directional 
arrangements of neighbours can therefore be insufficiently distinguished. This motivates 
the modified embedded-atom method (MEAM), which incorporates angular information into the 
effective density.

\subsection{Modified embedded-atom method}

The modified embedded-atom method, introduced by Baskes, extends the EAM formalism by 
incorporating angular contributions to the local atomic environment \cite{baskes1992meam}. 
The total energy retains the general embedding-plus-pair structure:

\begin{equation}
    E_{\mathrm{MEAM}}
    = \sum_i \left[
        F_{\alpha_i}(\bar{\rho}_i)
        + \frac{1}{2}\sum_{j\ne i}
        \phi_{\alpha_i\alpha_j}(r_{ij})
    \right],
    \label{eq:meam-energy}
\end{equation}

where $\bar{\rho}_i$ is the effective background density in the MEAM formalism. The distinction 
from conventional EAM is that this density incorporates information about the angular distribution 
of neighbouring atoms, rather than depending only on an isotropic scalar sum.

The MEAM energy thus combines two contributions. The embedding term describes the 
environment-dependent cohesive energy, while the pair term provides an explicit radial 
interaction between atoms. Both terms are parameterised, and their combined contribution 
determines the energy landscape sampled during structural relaxation and molecular dynamics.

\subsubsection{Background electron density and angular contributions}

In MEAM, the effective density is constructed from a spherically symmetric density contribution 
and angular corrections. A common representation is

\begin{equation}
    \bar{\rho}_i
    = \rho_i^{(0)} G(\Gamma_i),
    \label{eq:meam-effective-density}
\end{equation}

where $\rho_i^{(0)}$ is the zeroth-order, spherically averaged density, $G$ is a model-dependent 
screening or angular-density function, and $\Gamma_i$ collects contributions associated with the 
angular distribution of neighbouring atoms.

In the conventional formulation, the angular correction can be expressed schematically as

\begin{equation}
    \Gamma_i =
    \sum_{k=1}^{3} t_i^{(k)}
    \left(
        \frac{\rho_i^{(k)}}{\rho_i^{(0)}}
    \right)^2,
    \label{eq:meam-angular-correction}
\end{equation}

where $\rho_i^{(k)}$ denotes an angular partial-density measure of order $k$, and $t_i^{(k)}$ controls 
its relative contribution. The exact definitions of these partial densities, density scaling, and the function 
$G$ depend on the particular MEAM formulation and implementation. Equation~\ref{eq:meam-angular-correction} is 
therefore a schematic representation of the conventional angular-density construction, rather than a complete 
specification of every implementation option.

This angular dependence is important when modelling structures with different coordination geometries. For example, 
atoms in FCC, BCC, and other crystal structures occupy different spatial arrangements even when their nearest-neighbour 
distances are comparable. Directional information also becomes relevant when the coordination environment is altered by 
surfaces, defects, strain, or local chemical disorder. MEAM provides a framework for representing these effects within 
a computationally tractable interatomic potential \cite{baskes1992meam}.

\subsection{Relevance to surface-energy calculations}

Surface-energy calculations provide a particularly useful test of the environmental dependence represented by MEAM. Creating a free surface 
reduces the coordination of surface atoms and changes the local distribution of neighbours. Atoms in subsurface layers can also experience 
altered environments, while relaxation changes interatomic distances and local geometry. In chemically disordered alloys, the identities of 
surface and subsurface atoms introduce an additional source of variation in local bonding.

For a slab containing $N$ atoms, a common potential-energy-based estimate of the surface energy is

\begin{equation}
    \gamma =
    \frac{E_{\mathrm{slab}}-N e_{\mathrm{bulk}}}{2A},
    \label{eq:surface-energy}
\end{equation}

where $E_{\mathrm{slab}}$ is the total energy of the slab, $e_{\mathrm{bulk}}$ is the reference bulk energy per atom at a consistent composition and 
under compatible conditions, and $A$ is the area of one exposed face. The factor of two applies to a slab with two equivalent surfaces. The resulting 
energy difference measures the excess energy associated with creating the two surfaces relative to the chosen bulk reference.